\chapter{The Galois-theoretic proof}
\label{chap:galois}

This proof is complete in \texttt{Quintic/Galois.lean}.

\section{Radical values are solvable by radicals}

\begin{lemma}
  \label{lem:value-solvable-by-rad}
  \lean{Quintic.value_mem_solvableByRad}
  \leanok
  \uses{def:value}
  If $\mathrm{Value}(e, y)$, then $y$ lies in the subfield of $L$ of elements
  solvable by radicals over $K$.
\end{lemma}

\begin{proof}
  \leanok
  Induction on $e$. Each constructor of $\mathrm{Value}$ matches a closure property
  of Mathlib's \texttt{solvableByRad}.
\end{proof}

\section{The universal splitting field}

Let $x_0, \dots, x_4$ be fresh indeterminates, $S = \mathbb{C}[x_0, \dots, x_4]$ and
$E = \mathbb{C}(x_0, \dots, x_4)$. Write $e_k$ for the $k$-th elementary symmetric polynomial
in the $x_j$.

\begin{lemma}[Algebraic independence of the elementary symmetric polynomials]
  \label{lem:phi-injective}
  \lean{Quintic.Generic.φ_injective}
  \leanok
  \uses{def:K}
  The $\mathbb{C}$-algebra map $\varphi : R \to S$, $a_i \mapsto e_{5-i}$, is injective.
  Hence it extends to a field embedding $K \to E$, making $E$ an extension of $K$.
\end{lemma}

\begin{proof}
  \leanok
  $\varphi$ is a renaming of variables followed by the map $a_i \mapsto e_{i+1}$,
  which is injective by the fundamental theorem of symmetric polynomials
  (Mathlib's \texttt{esymmAlgHom\_injective}).
\end{proof}

\begin{lemma}[The quintic splits over $E$]
  \label{lem:map-generic-quintic}
  \lean{Quintic.Generic.map_genericQuintic}
  \leanok
  \uses{lem:phi-injective, def:generic-quintic}
  Over $E$, $P(X) = \prod_{j} \bigl(X - (-x_j)\bigr)$.
\end{lemma}

\begin{proof}
  \leanok
  Vieta's formulas: $\prod_j (X + x_j) = \sum_k e_k X^{5-k}$.
\end{proof}

\section{\texorpdfstring{$S_5$}{S5} in the Galois group}

\begin{lemma}[Permutations fix $K$]
  \label{lem:aut-hom}
  \lean{Quintic.Generic.autHom, Quintic.Generic.autHom_root}
  \leanok
  \uses{lem:phi-injective}
  Each permutation $\sigma$ of the $x_j$ induces an automorphism of $E$ over $K$,
  sending the root $-x_j$ to $-x_{\sigma(j)}$.
\end{lemma}

\begin{proof}
  \leanok
  Permuting variables fixes each $e_k$, so it fixes the image of $R$; extend to fraction
  fields with \texttt{IsFractionRing.fieldEquivOfAlgEquivHom}.
\end{proof}

\begin{lemma}[$S_5$ embeds in the Galois group]
  \label{lem:gal-hom-injective}
  \lean{Quintic.Generic.galHom, Quintic.Generic.galHom_injective}
  \leanok
  \uses{lem:aut-hom, lem:map-generic-quintic}
  Restricting the automorphisms of \cref{lem:aut-hom} to a splitting field of $P$ gives an
  injective group homomorphism $S_5 \to \mathrm{Gal}(P)$.
\end{lemma}

\begin{proof}
  \leanok
  Compare actions on the roots $-x_j$, which are pairwise distinct.
\end{proof}

\begin{corollary}
  \label{cor:gal-not-solvable}
  \lean{Quintic.Generic.not_isSolvable_gal}
  \leanok
  \uses{lem:gal-hom-injective}
  $\mathrm{Gal}(P)$ is not solvable.
\end{corollary}

\begin{proof}
  \leanok
  Subgroups of solvable groups are solvable, and $S_5$ is not solvable
  (Mathlib's \texttt{Equiv.Perm.not\_isSolvable}).
\end{proof}

\section{Conclusion}

\begin{lemma}[Roots are not solvable by radicals]
  \label{lem:root-not-solvable}
  \lean{Quintic.root_not_mem_solvableByRad}
  \leanok
  \uses{lem:aut-hom, lem:map-generic-quintic, cor:gal-not-solvable}
  No root of $P$ in $L$ is solvable by radicals over $K$.
\end{lemma}

\begin{proof}
  \leanok
  Let $y$ be a root and $q$ its minimal polynomial, so $q \mid P$. Over $E$, $q$ divides a
  product of linear factors, so it has some root $-x_{j_0}$. Applying the automorphisms of
  \cref{lem:aut-hom}, every $-x_j$ is a root of $q$, so $\deg q \geq 5$ and $q = P$.
  By Mathlib's \texttt{isSolvable\_gal\_minpoly}, $\mathrm{Gal}(q)$ would be solvable if $y$
  were solvable by radicals, contradicting \cref{cor:gal-not-solvable}.
\end{proof}

\begin{theorem}[Abel--Ruffini, Galois proof]
  \label{thm:galois-main}
  \lean{Quintic.quintic_unsolvable_galois}
  \leanok
  \uses{def:quintic-unsolvable, lem:value-solvable-by-rad, lem:root-not-solvable}
  The quintic is unsolvable.
\end{theorem}

\begin{proof}
  \leanok
  Combine \cref{lem:value-solvable-by-rad} and \cref{lem:root-not-solvable}.
\end{proof}
